\section{Introduction}
\label{sec:intro}

Reading underpins participation in education, employment, and daily life, yet for many people it is far from effortless.
The World Health Organization estimates that at least 2.2 billion people live with a near or distance vision impairment \cite{who2026vision}; age-related macular degeneration (AMD) alone is projected to affect 288 million people by 2040 \cite{wong2014amd}, and reading difficulties such as dyslexia \cite{peterson2012dyslexia} add to the population for whom text demands more than the visual or cognitive system can comfortably supply.
Digital text offers something print cannot: its presentation can change while readers read.
Font size, spacing, line length, and contrast are adjustable on screen and measurably affect reading \cite{rello2016makeitbig, beymer2008eyetracking, beier2021letterspacing, dyson2004layout, rubin1989contrast}, and an eye tracker supplies the moment-to-moment signal needed to apply an adjustment exactly when a reader needs it \cite{rayner1998eyemovements}.
A system that couples the two can deliver individualised typographic micro-interventions \cite{persson2025microintervention} without the reader ever issuing a command.

The promise carries a structural risk.
A typographic change that helps in principle also reflows the document, so the line a participant was reading moves, and an ill-timed change costs the reader their place in the very flow the system is trying to protect.
Prior work has begun to quantify this cost: interventions applied during reading differ in how much they disrupt it and how readily readers accept them \cite{pereira2024typography}, and highlighting the text read last before a change, combined with a gradual transition, gave the fastest resumption \cite{jensen2025context}.

We argue that the bottleneck in this research area has shifted.
Demonstrating that gaze can drive changes to the text being read no longer requires another prototype: such interfaces have been demonstrated repeatedly \cite{biedert2010eyebook, medan2024newline, rummens2026highlighting, jensen2025context}, including three adaptive readers built within one research programme in 2024 \cite{refsgaard2024rtr, pereira2024typography, palinec2024reactive}.
What is missing is the instrument for studying it: the programme's prototypes each wove sensing, decision logic, and the reading interface into one codebase, so each new study extended the previous system by forking it, and none offered a researcher console for operating a controlled session or a session record from which an experiment could be reconstructed.
On a coupled prototype, even a simple comparison such as fading in a font change gradually versus applying it abruptly means editing core code and maintaining two forks under identical sensing: an experiment that should be a configuration choice becomes a software project.

This paper presents a platform (Figure~\ref{fig:overview}) built to remove that friction.
A participant reads reflowable text on one screen while the text adapts to their gaze; on a second screen, a researcher follows the participant's gaze, fixations, and regressions live over a mirror of the same text, and applies typographic interventions or approves machine-proposed ones during the session.
Behind this workflow, the adaptive loop decomposes into four independently replaceable modules (sensing, eye-movement analysis, decision strategy, and intervention execution) connected through stable contracts, and writes every gaze sample, decision, and intervention into a schema-versioned session record that can be exported, replayed, and re-analysed.

The contributions are:
\begin{itemize}
  \item \textbf{A modular platform for gaze-contingent reading experiments}, in which sensing, analysis, decision, and intervention are separate modules behind stable contracts, extendable either inside the backend process by registering a module or from a separate process through the same contracts, carried over a WebSocket provider protocol.
  \item \textbf{A two-screen, researcher-in-the-loop workflow}, in which interventions are applied manually, proposed by a decision provider and gated by the researcher (advisory), or applied autonomously, without structural change between modes.
  \item \textbf{A context-preserving intervention mechanism that measures its own cost}: layout changes can be committed at controlled boundaries, a continuously captured anchor restores the reader's position across the reflow, and a graded positional outcome for every intervention is written into the session record.
  \item \textbf{A technical evaluation on real eye-tracking hardware} covering capture and transport, external-provider integration (including a collaborating study's reading-difficulty pipeline that analysed the evaluation sessions live), and context preservation, whose self-measurement exposed a defect in our original restore strategy that over-repositioned the reader on small reflows; we describe the revision and its geometric evidence.
\end{itemize}

We report the evaluation as an evaluation of the instrument: with four participants the behavioural observations are descriptive, and the paper makes no claim that any intervention improves reading.
The platform, its documentation, and the audit script that regenerates every quantity in this paper from the exported records are openly available.
